\cotefile{157-4}
\scanpage{348}
% @ head of sheet
\mspage{72}
% @ top, lines 1–4
une \inserted{(petite)} catégorie \unsure{filtrante} quelconque, et quitte à la \unsure{remplacer} par une catégorie \unsure{cofinale}, \unsure{pps|sps|p. ex.} $I$ un ensemble \unsure{ordonné}. L'axiome $\mathrm{Loc}(8\ \mathrm{bis})$ correspond au cas où $I=\mathbb{N}$.

% @ upper middle, lines 5–8
Notons que grâce à la construction de \unsure{la|l'é} \deleted{\unsure{équi|qu}} catégorie des chemins, nous avons une \deleted{\unsure{construction}} \inserted{factorisation}, fonctorielle en $f$, de \deleted{\unsure{toutes}} flèches $f\inserted{{:}\,X\to Y}$ de $\mathrm{Cat}$ en $f=pi$ :
% @ centre, first factorization
\diagram{\begin{tikzcd}[ampersand replacement=\&]
X \arrow[r,"i=i_f"] \& \overline X(f)=\overline X \arrow[r,"p=p_f"] \& Y{,}
\end{tikzcd}}
% @ middle, lines 10–12
avec $i\in W_{\unsure{\omega|W|w}}$, $p\in\mathrm{Parf}_{W_\omega}$, a fortiori $i\in W$, $p\in\mathrm{Parf}_W$. Ce foncteur $f\longrightarrow(\overline X(f),p_f)$ \unsure{commute} aux limites inductives.
% @ lower middle, lines 13–16
Ainsi, si on applique ceci à \unsure{chacun} des $f_\alpha:X_\alpha\longrightarrow Y_\alpha$, les factorisant en
\diagram{\begin{tikzcd}[ampersand replacement=\&]
X_\alpha \arrow[r,"i_\alpha",dash] \& \overline X_\alpha \arrow[r,"p_\alpha",dash] \& Y_\alpha{,}
\end{tikzcd}}
% @ foot, lines 17–21
\deleted{\unsure{on}} \unsure{on} trouve un système inductif de diagrammes de type $\Delta^2$ (factorisant les $f_i$), et en passant à la $\underrightarrow{\lim}$ on trouve
\diagram{\begin{tikzcd}[ampersand replacement=\&]
X \arrow[r,"i_{\deleted{\unsure{\alpha|\omega}}}"] \& \overline X \arrow[r,"p",dash] \& Y{,}
\end{tikzcd}}
qui n'est autre que la factorisation \unsure{canonique}
% archive sheet number: 347
